\subsection{EHR Clinical-Language and Clinical-Risk Prediction}
\label{sec:ehr}

We evaluate the unstructured and structured EHR modalities against classical floors under leakage-controlled protocols, and report an honest split in which the clinical transformer does \emph{not} dominate. Table~\ref{tab:ehr} summarizes all results.

\textbf{Clinical language.} The clinical-text corpus is MTSamples, a public-domain collection of de-identified transcribed medical reports, accessed through the Hugging Face mirror \texttt{rungalileo/medical\_transcription\_40} ($4{,}499$ notes)~. All results use grouped $5$-fold cross-validation in which each note is its own cross-validation group. We contrast a frozen clinical transformer, \texttt{emilyalsentzer/Bio\_ClinicalBERT}~\cite{clinicalbert}, used as a label-free embedding extractor, against a TF-IDF linear baseline; to avoid leakage the foundation-model embeddings are computed once with no label access, while the linear heads and the TF-IDF vocabulary are fit on training folds only. On Task~1, the binary surgery-vs-rest discrimination, Bio\_ClinicalBERT with a logistic-regression head attains AUROC $0.910$, clearly outperforming TF-IDF with logistic regression at $0.823$ and a hashing-plus-gradient-boosting configuration at $0.789$ (majority baseline $0.500$)~. On Task~2, the $40$-way specialty classification, the ordering reverses: TF-IDF with logistic regression reaches macro-AUROC $0.961$ (macro-F1 $0.437$), beating Bio\_ClinicalBERT with logistic regression at macro-AUROC $0.931$ (macro-F1 $0.405$)~. We report this split plainly: the frozen clinical transformer wins the coarse binary task but the classical sparse-lexical representation wins the fine-grained many-way task, so the transformer does not win everywhere and pretrained clinical embeddings are not uniformly superior to a well-tuned lexical baseline.

\textbf{Clinical risk (mortality).} For structured EHR risk we predict in-hospital mortality on the MIMIC-IV demo cohort using a gradient-boosted model over laboratory features, under patient-grouped $5$-fold cross-validation. The model attains AUROC $0.813$~. This estimate rests on a very small event count ($15$ positive events out of $252$ patients) and must therefore be read as indicative rather than validated; the wide implied confidence interval precludes any clinical claim.

\textbf{Diabetes complication risk (honest gap).} We deliberately report no model for diabetes complication risk. No open dataset carries real complication labels alongside the neural, wearable, and glucose signals used elsewhere in this work, so the target is not fine-tunable with honest supervision~. Rather than fabricate a label or a metric, the serving path abstains on this task. Consistent with the overall framing, these components constitute research-grade screening and decision support, not a validated diagnostic device.